\section{A direct quartic with 495 nonnegative unknowns}
\label{sec:projected-direct}

The direct branch encodes finite derivations in the repository's calibrated
numerical calculus $C_{222}$.  Its immediate parent is the incoming
synchronized-certificate construction with 522 nonnegative unknowns.
Combining its arithmetic certificates with branchwise projection of rule
labels gives the following fully materialized equation.

\begin{center}
\begin{tabular}{lrr}
\hline
Measure & Synchronized parent & Projected equation\\
\hline
Nonnegative unknowns &522&495\\
Quadratic residuals &400&391\\
Exact total degree &4&4\\
Expanded monomials &20,225&15,377\\
Maximum absolute coefficient &83,900,628&83,867,632\\
\hline
\end{tabular}
\end{center}

The actual polynomial is
\[
 P_{495}(z_0,\ldots,z_{494})=\sum_{j=1}^{391}R_j(z_0,\ldots,z_{494})^2.
\]
Every coefficient of every quadratic $R_j$, and every collected coefficient
of $P_{495}$, is supplied as an ordinary integer in the literal artifacts.
The notation does not quantify additional registers.  The source equations
below explain the construction; all their powers, binomials, factorials,
and bounded quantifiers have been removed from the delivered quartic.

This is a measured improvement of the direct, modest-coefficient quartic
branch.  It is not a claim that 495 is an optimal variable count.  The
construction does not encode a run of K215 and does not use Jones's
universal polynomial.  Its ordinary consistency theorem imports the
identified $C_{222}$ calibration; it is not a new end-to-end proof-assistant
certificate.

\subsection{Project the premise and conclusion labels first}

Write
\[
 I(a,b)=2(a+b)(a+b+1)+4b,\qquad A(a,b)=I(a,b)+2.
\]
The calibrated source has twelve branches.  Its truth and axiom branches
have premise labels $p=q=0$ and conclusion label
$v=h_b(a,b,c,d)$, with truth represented by $h_b=0$.
The exact nine nonzero conclusion polynomials are the inherited numerical
constructors.  The two inference branches admit the parameterizations
\[
 (p,q,v)=(a,I(a,b),b),\qquad
 (p,q,v)=(b,0,A(a,b)).
\]
For modus ponens these parameters are the old first premise and conclusion;
for generalization they are its generalization parameter and premise.
Every old instance therefore has a new instance, and substitution recovers
an old instance from every new one.  Since existential quantification
distributes over a finite disjunction, unused parameters in an unselected
branch impose no compatibility condition.

There is one further local saving.  The bundle branch's conclusion is
independent of its fourth parameter.  Its guard
\[
 g^2=(d_{\rm old}+1)^2,
 \qquad g=(c-a)(c-b),\qquad d_{\rm old}\ge0,
\]
can be replaced by $g^2=d+1$, $d\ge0$.  Both existential conditions mean
$g\ne0$.  An old witness gives $d=d_{\rm old}^2+2d_{\rm old}$; a new
witness gives $d_{\rm old}=|g|-1$.  The fourth parameter in another
selected branch retains its original meaning.

Fix outer parameters $t,m\ge1$ and $x,y,M_0\ge0$.  Set $M=M_0+1$ and
$w(j)=1+Mj$.  For each branch let
\[
 e_b=S' w(v_b)-S w(p_b)w(q_b).
\]
Let $e_0,\ldots,e_8$ be the truth and eight unguarded axiom residuals;
let $e_{\rm MP},e_{\rm GEN},e_{\rm B}$ be the two inference and bundle
residuals.  Put $G=((c-a)(c-b))^2-d-1$.  The complete row equation is
\begin{align*}
0={}&\left(e_0\cdots e_8e_{\rm MP}e_{\rm GEN}\right)^2
                   (e_{\rm B}^2+G^2)\\
 &+\bigl[x-S-(1+(i+1)y)u\bigr]^2
  +\bigl[x-S'-(1+(i+2)y)v\bigr]^2.
\end{align*}
Every summand is nonnegative.  Consequently this equation is equivalent
to the two beta reconstructions together with one valid selected rule's
accumulator identity.  Only eight bounded row coordinates remain:
\[
 (a,b,c,d,S,S',u,v),\qquad
 (a,b,c,d,S,S',u,v)\le(m,m,m,m,y,y,x,x)
\]
coordinatewise.  The implementation also shares $S'M$ in
$S'w(h)=S'+(S'M)h$ and squares the product of the eleven unguarded
residuals once.  These are exact polynomial identities, not probabilistic
fingerprint equalities.

\subsection{A small circuit bounds all labels}

At $a=b=c=d=m$, a literal certificate bounds every constructor by a fixed
power of $R_0=m+2$.  Assign exponent one to $m$, exponent
$\lceil\log_2 c\rceil$ to a positive numeral $c$, exponent zero to zero,
and propagate
\[
 e_{U+V}=\max(e_U,e_V)+1,\qquad e_{UV}=e_U+e_V.
\]
Since $R_0\ge2$, induction proves $U\le R_0^{e_U}$ at every gate.
All constructor coefficients are nonnegative, so the diagonal evaluation
also bounds the labels when the four parameters are merely at most $m$.
The largest certified exponent is $3,445,903$; use
\[
 h(m)=(m+2)^{4,194,304},\qquad M=(h(m)+1)^t.
\]
The fixed power defining $h$ takes 22 squarings.  This high semantic degree
does not raise the quartic degree of its arithmetic graph.

Here is the exact carrier argument.  For the $t$ selected rule instances
define
\[
 B_0(X)=\prod_{i<t}(1+v_iX),\qquad
 B_1(X)=(1+X)\prod_{i<t}(1+p_iX)(1+q_iX).
\]
The endpoints $S_0=M+1$, $S_t=1$ and the accumulator identities imply
$B_0(M)=B_1(M)$.  All states are positive, by induction from the first
endpoint and positivity of the label weights.  With $h=h(m)$,
\[
 B_0(M)\le(1+hM)^t<((h+1)M)^t=M^{t+1}.
\]
Thus $\deg B_1\le t$, since a nonnegative integer polynomial of greater
degree would have value at least $M^{t+1}$.  Both polynomials now have
coefficient sum at most $(h+1)^t=M$, and constant coefficient one.  Each
coefficient is therefore strictly smaller than $M$.  Uniqueness of
base-$M$ digits proves $B_0=B_1$.  Comparing the multiplicities of the
distinct factors $1+jX$, $j>0$, gives the exact Horn balance for target
label one.  The inherited closure-separation argument turns that balance
into a finite derivation, without excluding cycles or repeated premises.

Conversely, expand a finite derivation to a finite tree.  Choose $t$ for
its rule occurrences and choose $m$ afterward to bound their four
parameters, including each new bundle slack $g^2-1$.  Products of the
weights of outstanding obligations give the required integral
accumulator trace.  The independent parameter $m$ avoids any need to pad
the derivation until its horizon bounds its own labels.

\subsection{Compile and contract}

Choose $y$ above the finitely many states and divisible by $t!$, then
choose a beta code $x$ for the trace.  Its natural quotients are at most
$x$.  Conversely, the bounds $S,S'\le y$ make both reconstructions
canonical remainders, since $y<1+(i+1)y$.

A second literal majorant certificate uses base $2+t+x+y+m+M_0$ and
fixed exponent $8,388,608$.  It bounds both nonnegative sides of the
row equation at all upper row bounds, as well as the index and the
individual bounds.  The inherited fixed-vector CRT theorem then applies
to the eight columns.  Its factorial factors cancel because the CRT
moduli are one modulo the chosen common factorial.

The resulting finite arithmetic circuit makes 11 binomial requests and
24 positive-power requests, with 19 distinct Pell indices and 10 distinct
power bases.  The shared binomial radix and common Pell-frame theorems
apply to this finite family unchanged.  In soundness one identifies the
power graphs first, the binomial digits second, and the factorial/CRT
semantics afterward.  In completeness one chooses the intended acyclic
semantic values first and the shared auxiliary certificates afterward.
This order avoids circular justification of the shared radix or frame.

The emitted parent of the final quartic has 823 natural variables and
719 quadratic residuals.  Exact nonnegative monic substitutions eliminate
328 further variables and residuals.  The inverse substitutions have
nonnegative integer coefficients, and the retained residuals remain
quadratic.  Thus both existence implications hold over the naturals.
The independent replay verifies all 328 substitutions and the complete
15,377-term sum-of-squares expansion.

\section{Recover nonnegativity from the equations themselves}
\label{sec:frame-domain}

Converting every natural coordinate independently would waste many
variables.  The displayed equations already force most coordinates to
have the required signs.  A finite certificate shows that guarding only
170 of the 495 coordinates suffices.  The main additional argument uses
the common Pell frame before invoking any exponentiation semantics.

\subsection{Local sign implications}

Suppose a literal residual can be written
\[
 A(\mathbf z)z_i=B(\mathbf z),
\]
where neither $A$ nor $B$ involves $z_i$.  If previously established facts
prove $A>0$ and $B\ge0$, then $z_i\ge0$; if $B>0$, then $z_i>0$.
The certificate uses this elementary implication and the weakening
$z_i>0\Rightarrow z_i\ge0$.  Its sign checks allow even powers without
assuming a sign for their bases.  Positive integer monomials can also
pay an explicitly bounded negative constant term.

Starting with the 170 declared guards, 531 such checked implications
prove 447 of the 495 coordinates nonnegative.  Crucially, this first
stage proves all prerequisites of the following common-frame argument.
No $T$ or Pell abscissa whose sign is recovered by that argument is used
prematurely.

\subsection{One frame proves 29 further signs}

The exact raw compiler identities restore polynomially from the final
495 coordinates.  Let $W,A,Z$ be its common frame, and let
$(n,k,m)$ range over the finite power requests.  The frame identities are
\[
 W=2+\sum_{\xi\in\mathcal A}\xi,\qquad
 A^2=1+W^3(W+2)Z^2.
\]
Here $\mathcal A$ is the finite set of distinct base, exponent, and output
atoms in the circuit; distinct atoms with equal values are still counted
separately.  The exact list is supplied.  Before the new argument,
all these atoms are nonnegative, every base $n$ is strictly positive,
$A\ge0$, and $Z\ge1$.  These facts are checked through the literal
restorations, including the factorial output used as a power base.
Consequently
\[
 W\ge2,\qquad 1\le n\le W,\qquad A>W^2.
\]
The last strict inequality follows immediately by squaring; no Pell
classification or already-established power semantics is needed.

Each of the ten distinct base moduli satisfies
\[
 T_n=2An-n^2-1.
\]
Therefore $T_n\ge2A-W^2-1>0$.  For each power request the exact oriented
congruence is
\[
 X+Yn=YA+m+T_nq,
\]
and the first sign stage has already established $Y,m,q\ge0$.
Since $A-n>0$, rearranging gives
\[
 X=Y(A-n)+m+T_nq\ge0.
\]
This recovers the signs of all 19 distinct index abscissae.  After adding
the ten modulus signs and these 19 abscissa signs, 21 further local
implications prove all 495 coordinates nonnegative.

The order matters: base positivity precedes the frame argument, modulus
positivity precedes abscissa recovery, and full natural-domain arithmetic
semantics is invoked only after every coordinate has been recovered.
This separates a short ordered-ring proof from the much stronger
Pell-index argument used later to identify actual powers.

\subsection{A quartic in 977 unrestricted integers}

For an integer $x$, the classical three-square theorem gives
\[
 x\ge0\quad\Longleftrightarrow\quad
 \exists u,v,w\in\mathbb Z:\quad x=u^2+v^2+w(w+1).
\]
Soundness is immediate.  For completeness, $4x+1$ has a representation
by three squares; modulo four exactly one square is odd.  Writing that
odd integer as $2w+1$ and the even integers as $2u,2v$ proves the formula.

Adjoin this quadratic identity for each of the 170 guarded coordinates.
The other coordinates require no separate representation, because the
checked domain-recovery theorem forces their nonnegativity.  Keeping the
original coordinates at this stage preserves degree two in every
residual.  The intermediate unrestricted-integer system therefore has
\[
 495+3\cdot170=1005\quad\text{variables},\qquad
 391+170=561\quad\text{quadratic residuals}.
\]
Twenty-eight further signed monic substitutions preserve quadratic
degree and existence over the integers.  The final literal equation has
\[
 \boxed{\begin{gathered}
 977\ \text{unrestricted integer unknowns},\qquad
 533\ \text{quadratic residuals},\qquad \deg P_{977}=4,\\
 24,632\ \text{expanded monomials},\qquad
 \max|\text{coefficient}|=83,869,250.
 \end{gathered}}
\]
The comparison with the synchronized parent's 2088-integer quartic is
a reduction of 1111 integer variables, or approximately $53.2\%$, while
retaining degree four.  This is a comparison within the direct branch;
the repository's much smaller-variable high-degree constructions remain
different tradeoffs.

Every integer zero of $P_{977}$ restores the guarded 495-coordinate
system, hence a completely nonnegative zero of $P_{495}$ and a calibrated
finite refutation.  Conversely, a refutation gives a natural zero,
three-square witnesses for the 170 guards, and then a zero of the
contracted integer system.  Composing with the named ordinary source
calibration and negating yields the desired ZFC-provable biconditional
between absence of integer zeros and $\operatorname{Con}(\mathrm{ZFC})$.

\subsection{What the finite verifier establishes}

The domain checker is independent of the sign-discovery code.  It
rederives every local implication from a literal quadratic residual and
checks 1535 raw polynomial restorations.  Of the raw equations, 1040
restore to zero identities and 391 restore to final residuals.  It also
binds all 1243 raw arithmetic-gate metadata records to their exact
defining residuals before expanding them.  Finally it checks the frame
norm and sum, ten modulus identities, and all 24 oriented power
identities.  The ordinary inequalities above then justify its 29
additional sign conclusions.

Corruption controls require rejection after a gate operand is changed,
a guard is deleted, a pivot orientation is reversed, a frame conclusion
is omitted, or a modulus identity is replaced.  Normal and optimized
Python runs perform the checks explicitly.  These finite verifications
certify the formulas, transformations, and coefficient data; they do not
replace the unbounded source proof or decide whether a refutation exists.
